\subsection{Ley de Amdahl}
\label{section:amdahl}

La \emph{ley de Amdahl} establece, a grandes rasgos, que la aceleración que se puede tener en un programa paralelo está limitado por la fracción secuencial (no paralelizable) del programa, independientemente del número de procesadores o núcleos que se añadan.

Teniendo en cuenta lo anterior, entonces el tiempo de ejecución paralelo $T_p$ se puede definir como la suma del tiempo de ejecución de la fracción secuencial $f$ y del tiempo de ejecución de la fracción paralela $\left(1 - f\right)$, esto es:

\[T_p = f \cdot T_s + \left(1 - f\right) \frac{T_s}{p}\]

En consecuencia, podemos redefinir la aceleración $S$ como:

\[S = \frac{T_s}{f \cdot T_s + \left(1 - f\right) \cdot T_s / p} = \frac{1}{f + \left(1 - f\right) / p}\]

Por lo tanto, de acuerdo con la ley de Amdahl, independientemente del número de procesadores o núcleos que se añadan ($p$ se hace cada vez más grande y $\left(1 - f\right) / p$ se hace cada vez más pequeño), la aceleración máxima que se puede tener es:

\[S \le \frac{1}{f}\]

Por ejemplo, si se considera que la parte secuencial de un programa abarca $f = 10\%$ de todo el programa paralelo, entonces a lo mucho tendríamos una aceleración de $10$, sin importar si tenemos miles de procesadores o núcleos. \cite{pacheco2011introduction}